% Dibangkitkan oleh revisi/hasil/gerbang_batch/bab5_jalan.py. Jangan diubah dengan tangan.
\begin{table}[htbp]
\caption{Analisis kepekaan: 35 instans benchmark saja (selang 95\%)}
\label{tab:uji-kepekaan}
\centering\footnotesize
\setlength{\tabcolsep}{4pt}
\begin{adjustbox}{max width=\textwidth}
\begin{tabular}{llll}
\toprule
Pembanding & $\lambda$ & Uk.~1 (\%) & Uk.~2 (\%) \\
\midrule
Random + penjaga & 0{,}3 & $-23{,}96$ [$-26{,}09$; $-21{,}73$]$^{*}$ & $-23{,}18$ [$-24{,}80$; $-21{,}60$]$^{*}$ \\
Random + penjaga & 3 & $-15{,}06$ [$-16{,}55$; $-13{,}55$]$^{*}$ & $-15{,}51$ [$-16{,}66$; $-14{,}15$]$^{*}$ \\
Fixed-L3 + penjaga & 0{,}3 & $+17{,}47$ [$+8{,}03$; $+29{,}39$]$^{*}$ & $-5{,}02$ [$-6{,}36$; $-2{,}21$]$^{*}$ \\
Fixed-L3 + penjaga & 3 & $+0{,}22$ [$-2{,}17$; $+3{,}19$] & $-6{,}38$ [$-7{,}40$; $-4{,}70$]$^{*}$ \\
\bottomrule
\end{tabular}
\end{adjustbox}
\par\smallskip\footnotesize\textit{Catatan:} Angka adalah selisih $J$ ARLS-PPO terhadap pembanding dalam persen; negatif berarti ARLS-PPO lebih baik. Uk.~1 rata-rata rasio per kondisi, uk.~2 rasio $J$ total; setiap instans berbobot sama. Tanda $^{*}$: selang tidak memuat nol.
\end{table}
